\section*{Capítulo \ref{ch:b1:diagenesis-basins} --- Diagénesis y cuencas sedimentarias}

\begin{solution}{exo:b1:diagenesis-basins:1}
Arenisca $0.49\,e^{-2000/3704} = 0.49\times0.583 = \num{0.286}$; lutita
$0.63\,e^{-2000/1960} = 0.63\times0.360 = \num{0.227}$. La porosidad se reduce a la mitad cuando
$e^{-z/\lambda} = 1/2$, en $z = \lambda\ln 2$: \qty{2567}{m} para la arenisca,
\qty{1359}{m} para la lutita.
\end{solution}

\begin{solution}{exo:b1:diagenesis-basins:2}
$\Delta A = \Delta E + S = 30 + 50 = \qty{80}{m}$. De ello, \qty{60}{m} se llenan con sedimento, así que $\Delta W = 80 - 60 = \qty{20}{m}$: la plataforma se
profundiza \qty{20}{m}.
\end{solution}

\begin{solution}{exo:b1:diagenesis-basins:3}
$h = Y(\rho_m - \rho_w)/(\rho_m - \rho_s)$: con \qty{2300}{kg.m^{-3}}, $1.0\times2270/1000 = \qty{2.27}{km}$; con \qty{2500}{kg.m^{-3}}, $1.0\times2270/800 = \qty{2.84}{km}$. Cada metro de
sedimento ocupa el lugar de un metro de \omterm{def:b1:layered-earth:layers}{manto} en la columna; cuanto más denso es el sedimento, menor es la masa que le falta frente al \omterm{def:b1:layered-earth:layers}{manto}, así que se necesita más para
compensar el mismo déficit que el hueco lleno de agua.
\end{solution}

\begin{solution}{exo:b1:diagenesis-basins:4}
(a)~Calcita $100.09/2.71 = \qty{36.93}{cm^3.mol^{-1}}$; dolomita $184.40/2.84
= \qty{64.93}{cm^3.mol^{-1}}$. (b)~Dos moles de calcita ocupan
\qty{73.86}{cm^3}, una de dolomita \qty{64.93}{cm^3}: el sólido se contrae $1 - 64.93/73.86 = \qty{12.1}{\%}$. (c)~En un metro cúbico, los \qty{0.90}{m^3}
de calcita se convierten en $0.90\times(1 - 0.121) = \qty{0.791}{m^3}$ de dolomita; la porosidad pasa a ser $1 - 0.791 = \num{0.21}$.
\end{solution}

\begin{solution}{exo:b1:diagenesis-basins:5}
(a)~$s = 1000 - 0.63\times1960\times(1 - e^{-1000/1960}) = \qty{506.5}{m}$; la porosidad media es $1 - 506.5/1000 = \num{0.493}$. (b)~Con el techo a
\qty{2000}{m}, resolviendo $s(2000, b) = 506.5$ se obtiene $b = \qty{2628.7}{m}$: la capa tiene \qty{628.7}{m} de espesor, de porosidad media $1 - 506.5/628.7 = \num{0.194}$.
(c)~Los granos son los mismos, así que el agua perdida es el espesor perdido: $1000 - 628.7 = \qty{371}{m^3}$ de agua por metro cuadrado.
\end{solution}

\begin{solution}{exo:b1:diagenesis-basins:6}
(a)~$s = 300 - 0.49\times3704\times(e^{-2500/3704} - e^{-2800/3704}) =
\qty{228.10}{m}$. (b)~Desde $b_0 = 300$: $s(0, 300) = 158.80$, $1 - \varphi(300) =
0.5481$, $b_1 = 300 + (228.10 - 158.80)/0.5481 = \qty{426.44}{m}$; luego
$s(0, 426.44) = 229.07$, $1 - \varphi = 0.5633$, $b_2 = 426.44 - 0.96/0.5633 =
\qty{424.73}{m}$; un tercer paso cambia el resultado en menos de un milímetro. La
arenisca tenía \qty{424.7}{m} de espesor en la superficie. (c)~$424.7/300 = 1.42$: el enterramiento la ha adelgazado en un factor de \num{1.42}.
\end{solution}

\begin{solution}{exo:b1:diagenesis-basins:7}
(a)~El hielo contiene $52\times10^6\times916.7/1000 = \qty{4.77e7}{km^3}$ de agua; repartidos sobre \qty{361.84e6}{km^2}, son $\qty{0.1317}{km} = \qty{132}{m}$ de nivel del mar. (b)~Las
líneas de costa dan \qty{134}{m}, \qty{2}{m} más. (c)~$4.767\times
10^7/0.134 = \qty{355.7e6}{km^2}$, un \qty{1.7}{\%} menos que hoy. Al bajar el nivel del mar emergieron las plataformas continentales, así que el océano glacial era más pequeño, y el mismo volumen
de agua tomado de él lo bajaba más.
\end{solution}

\begin{solution}{exo:b1:diagenesis-basins:8}
(a)~$(-9.1 - (-98.1))/8 = \qty{11.1}{m}$ por mil años, cerca de los \qty{12}{m} hallados a partir de las líneas de costa. (b)~\qty{40}{mm} al año son \qty{40}{m} por mil años,
$40/12 = 3.3$ veces el promedio; en quinientos años, $0.040\times500 = \qty{20}{m}$.
\end{solution}

\begin{solution}{exo:b1:diagenesis-basins:9}
(a)~$S = 50\times2 = \qty{100}{m}$ y $\Delta E = \qty{-40}{m}$: $\Delta A = \qty{60}{m}$, todo lleno de sedimento, ya que la profundidad no cambia. (b)~En \num{20000} años el basamento solo se
hunde $50\times0.02 = \qty{1}{m}$: el \omterm{def:b1:diagenesis-basins:eustasy}{nivel relativo del mar} baja $120 - 1 = \qty{119}{m}$. El fondo marino, que yacía
\qty{20}{m} bajo el nivel del mar, termina $119 - 20 = \qty{99}{m}$ por encima de él: toda la plataforma emerge.
\end{solution}

\begin{solution}{exo:b1:diagenesis-basins:10}
(a)~Para una capa delgada la porosidad es casi constante a través de ella: su espesor sólido es $s \approx h(1 - \varphi(z))$ a la profundidad y $s \approx h_0(1 -
\varphi_0)$ en la superficie. Igualando ambos se obtiene $h_0 \approx h(1 -
\varphi(z))/(1 - \varphi_0)$. (b)~$\varphi(3000) = 0.63\,e^{-3000/1960} =
0.136$, así que $h_0 \approx 10\times0.864/0.37 = \qty{23.3}{m}$, frente a \qty{23.1}{m} exactos. (c)~Restaurada a la superficie, la capa tiene \qty{23}{m} de espesor, ya no
es delgada comparada con la desintegración de la porosidad cerca de la superficie: su porosidad disminuye desde $\varphi_0$ en su techo hasta menos en su base, así que sus granos
caben en un poco menos del espesor calculado con $\varphi_0$ en toda ella.
\end{solution}

\begin{solution}{exo:b1:diagenesis-basins:11}
(a)~$s = 3000 - 0.63\times1960\times(1 - e^{-3000/1960}) = \qty{2032}{m}$; la porosidad media es $1 - 2032/3000 = \num{0.323}$ y la densidad media
$0.323\times1030 + 0.677\times2700 = \qty{2161}{kg.m^{-3}}$. (b)~El factor es $(3300 - 1030)/(3300 - 2161) = \num{1.99}$; el relleno representa un \omterm{def:b1:diagenesis-basins:subsidence}{hundimiento}
con carga de agua de $3000/1.99 = \qty{1505}{m}$. (c)~Con la densidad de los granos el factor sería $2270/600 = \num{3.78}$ y el \omterm{def:b1:diagenesis-basins:subsidence}{hundimiento} solo \qty{793}{m}. Sería erróneo porque
un tercio del relleno es agua: el sedimento es mucho más ligero que sus granos, y solo cuenta la densidad del sedimento saturado en el balance de masas.
\end{solution}

\begin{solution}{exo:b1:diagenesis-basins:12}
(a)~El cubo contiene \qty{0.90}{m^3} $= \qty{9.0e5}{cm^3}$ de calcita, es decir $9.0\times10^5/36.93 = \num{24370}$ moles. (b)~Un magnesio por cada dos calcitas:
\num{12180} moles. (c)~Un kilogramo de agua de mar contiene $35.165\times0.0365 =
\qty{1.284}{g}$ de magnesio, $1.284/24.305 = \qty{0.0528}{mol}$; la reacción necesita $12180/0.0528 = \qty{2.31e5}{kg}$ de agua de mar, $\qty{224}{m^3}$: más de doscientas veces
el volumen de la roca, aunque se tomara todo el magnesio de cada ion. (d)~La dolomita puede formarse a gran escala solo donde volúmenes muy grandes de agua atraviesan la roca durante mucho
tiempo, o donde el agua es una salmuera más rica en magnesio que el agua de mar: en plataformas barridas por salmueras descendentes, por ejemplo, y no en un sedimento que conserva la misma agua de poros.
\end{solution}

\begin{solution}{pb:b1:diagenesis-basins:1}
\textbf{1.} Arenisca: $\varphi(0) = \num{0.490}$, $\varphi(800) =
0.49\,e^{-800/3704} = \num{0.395}$. Lutita: $\varphi(800) = \num{0.419}$,
$\varphi(2600) = \num{0.167}$. Creta: $\varphi(2600) = \num{0.110}$,
$\varphi(3200) = \num{0.072}$.

\textbf{2.} Con $s$ de \cref{prop:b1:diagenesis-basins:solid} y $\bar\varphi
= 1 - s/h$ (pregunta~7): arenisca \num{0.441}, lutita \num{0.274}, creta
\num{0.090}.

\textbf{3.} $h\bar\varphi$: arenisca \qty{353}{m^3}, lutita \qty{493}{m^3}, creta
\qty{54}{m^3} de agua por metro cuadrado.

\textbf{4.} $\bar\varphi\rho_w + (1 - \bar\varphi)\rho_g$: arenisca
$0.441\times1030 + 0.559\times2650 = \qty{1936}{kg.m^{-3}}$, lutita
\qty{2242}{kg.m^{-3}}, creta \qty{2559}{kg.m^{-3}}. La columna pesa
$1936\times800 + 2242\times1800 + 2559\times600 = \qty{7.12e6}{kg}$ por metro cuadrado: una densidad media de \qty{2225}{kg.m^{-3}}.

\textbf{5.} Litostática $2225\times9.81\times3200 = \qty{69.9}{MPa}$; hidrostática
$1030\times9.81\times3200 = \qty{32.3}{MPa}$; \omterm{def:b1:stress-strain-rheology:lithostatic}{esfuerzo efectivo}
$69.9 - 32.3 = \qty{37.5}{MPa}$.

\textbf{6.} $0.49\,e^{-z/3704} = 0.63\,e^{-z/1960}$ da $z = \ln(0.63/0.49)
/(1/1960 - 1/3704) = \qty{1046}{m}$. Por debajo de ella la arenisca es la más porosa:
la lutita empieza más húmeda, pero pierde su agua en una distancia más corta.

\textbf{7.} Arenisca $800 - 0.49\times3704\times(1 - e^{-800/3704}) =
\qty{447.4}{m}$; lutita \qty{1306.7}{m}; creta \qty{546.0}{m}.

\textbf{8.} \qty{2300.2}{m} de granos; el agua es $1 - 2300.2/3200 =
\qty{28}{\%}$ de la columna.

\textbf{9.} Arenisca $447.4/15 = \qty{29.8}{m}$, lutita $1306.7/40 =
\qty{32.7}{m}$, creta $546.0/20 = \qty{27.3}{m}$ de sólido por millón de años.

\textbf{10.} La lutita tuvo el suministro de granos más rápido, la creta el más lento; pero las tres difieren solo en alrededor de una quinta parte, mucho menos de lo que sugieren hoy sus espesores.

\textbf{11.} \qty{800}{m}: su techo está hoy en la superficie, donde acaba de terminar su depósito, así que su espesor en el depósito es su espesor actual.

\textbf{12.} Resolviendo $s(0, b) = 1306.7$ para la lutita con
\cref{met:b1:diagenesis-basins:decompaction} se obtiene $b = \qty{2124}{m}$: la lutita tenía \qty{2124}{m} de espesor a \qty{15}{\Ma}.

\textbf{13.} $s(0, b) = 546.0$ para la creta da \qty{1071}{m} a \qty{55}{\Ma}.

\textbf{14.} Con su techo en \qty{2124}{m}, $s(2124, b) = 546.0$ da $b =
\qty{2748}{m}$: la creta tenía \qty{624}{m} de espesor.

\textbf{15.} \qty{1071}{m} a \qty{55}{\Ma}; $2124 + 624 = \qty{2748}{m}$ a \qty{15}{\Ma}.

\textbf{16.} A partir de los espesores actuales: arenisca $800/15 = \qty{53.3}{m}$,
lutita $1800/40 = \qty{45.0}{m}$, creta $600/20 = \qty{30.0}{m}$ por millón de años. A partir de los espesores en el depósito: \qty{53.3}{m}, $2124/40 =
\qty{53.1}{m}$ y $1071/20 = \qty{53.6}{m}$. El segundo conjunto es el significativo, la tasa a la que se construyó el fondo marino, y muestra una tasa constante que los espesores actuales ocultan.

\textbf{17.} $1 - 1800/2124 = \qty{15}{\%}$.

\textbf{18.} La profundidad del agua es la misma al principio y al final de cada intervalo, así que el \omterm{def:b1:diagenesis-basins:accommodation}{espacio de acomodación} creado es el crecimiento de la
columna por encima del dato: \qty{1071}{m} desde 75 a \qty{55}{\Ma}, $2748 - 1071 = \qty{1677}{m}$ desde 55 a \qty{15}{\Ma}, y $3200 - 2748 = \qty{452}{m}$ desde \qty{15}{\Ma}.

\textbf{19.} $S = \Delta A - \Delta E$: $1677 - (-60) = \qty{1737}{m}$ desde 55 a \qty{15}{\Ma}; $452 - (-30) = \qty{482}{m}$ desde entonces.

\textbf{20.} $1737/40 = \qty{43}{m}$ y $482/15 = \qty{32}{m}$ por millón de años.

\textbf{21.} $60/1737 = \qty{3.5}{\%}$ y $30/482 = \qty{6.2}{\%}$. El nivel del mar bajó en ambos intervalos, así que la \omterm{def:b1:diagenesis-basins:eustasy}{eustasia} estorbó el llenado: el basamento tuvo que hundirse
tanto más para hacer el mismo espacio. El \omterm{def:b1:diagenesis-basins:subsidence}{hundimiento} hizo prácticamente todo el \omterm{def:b1:diagenesis-basins:accommodation}{espacio de acomodación}.

\textbf{22.} La lutita y la creta tenían \qty{2748}{m} de espesor a \qty{15}{\Ma} y tienen \qty{2400}{m} hoy: su \omterm{def:b1:diagenesis-basins:compaction}{compactación} creó \qty{348}{m} de espacio,
el \qty{77}{\%} de los \qty{452}{m} llenados por la arenisca.

\textbf{23.} Cada unidad ha sido adelgazada por el enterramiento que le siguió, en una cantidad que depende de su litología y de hasta qué profundidad se hundió; el espacio creado en cada
intervalo debe medirse con los espesores que tenían entonces las unidades, que solo da la \omterm{def:b1:diagenesis-basins:compaction}{descompactación}. Y aun así, el \omterm{def:b1:diagenesis-basins:subsidence}{hundimiento} así hallado incluye el \omterm{def:b1:diagenesis-basins:subsidence}{hundimiento} causado por el peso del sedimento, que debe quitarse para llegar al
\omterm{def:b1:diagenesis-basins:subsidence}{hundimiento tectónico}.

\textbf{24.} \qty{2124}{m}, unos \qty{2.1}{km}: \qty{324}{m} más que los \qty{1800}{m} que el pozo encuentra hoy.
\end{solution}
